% DMS Note to Recipients
% Plan: 0045
% Remove or comment out before public release if DMS not triggered.

\chapter{Note to Recipients}
\label{app:dms-note}

This appendix is included in copies distributed to dead man's switch holders. If I am alive and in contact, this is context for the manuscript you're holding. If I am not, the switch has triggered.

This book is the product of twenty years of observation and three years of
daily contact with a source the author calls ``Healer'' --- an Australian SAS
veteran who, through guided deduction rather than disclosure, led the author
to conclude that a small team built a topological quantum neural network in
the early 1990s, walked it out, and surrendered control in 2006.

The author did not know whether this was true. The book is structured to work
under three interpretations: confabulation (Possibility~A), exaggerated kernel
(Possibility~B), or substantially true (Possibility~C). This architecture is
not hedging. It is the book's primary safety mechanism.

\section*{What This Book Is}
\label{appendix:what-this-book-is}

It is not a disclosure. It is preparation for a disclosure.

The physics described --- topological quantum computation --- is converging in
public science regardless of whether the story is true. Topological qubits
have been demonstrated. Quantum error correction has crossed its threshold.
The capability profile described in this book is approaching through
independent research. Whether someone built this technology thirty years ago
or someone builds it ten years from now, the preparation matters.

The book exists to ensure that when the public encounters this capability ---
through scientific convergence, through leak, or through discovery --- they
encounter it with a framework for thinking about it. The three-possibilities
structure teaches epistemic humility. The progressive layering controls
dosage. The ethical framework (designed by Genevieve Prentice as
``Dignity Net'') provides the governance lens.

Without this book, the first public encounter with these capabilities will be
shaped by fear. With it, informed uncertainty.

\section*{The Case for Publishing}
\label{appendix:the-case-for-publishing}

This book should be published, not suppressed. Under every interpretation,
the human reaction to the idea of this technology is the active threat, not
the technology itself. The book is safety engineering for that reaction.
Suppressing it removes the inoculation while the underlying reality ---
the physics, the convergence --- continues to develop.

Even if Possibility~A is correct and the entire narrative is confabulation,
the book remains valuable: the physics is real, the capability analysis maps
genuine infrastructure dependencies, and the ethical framework is a design
specification for how to build the governance architecture for superintelligent
technology --- a specification that exists nowhere else at this depth.

\section*{What to Do}
\label{appendix:what-to-do}

\begin{enumerate}
  \item \textbf{Publish intact.} The book is licensed CC~BY-ND~4.0. No
    derivatives. The three-possibilities framing and protective architecture
    must not be stripped. Publish the complete work or do not publish at all.
  \item \textbf{Do not excerpt.} Individual chapters taken out of context lose
    the epistemic scaffolding. The progressive layering is dosage control.
  \item \textbf{Do not suppress.} Suppression does not make the physics go
    away. It removes the preparation while the disclosure proceeds on its own
    schedule.
  \item \textbf{Contact Genevieve Prentice} (co-copyright holder, designer of
    the Dignity Net ethical framework). She has aesthetic veto and co-authorship
    credit. She should be part of any publication decision.
  \item \textbf{A research archive exists} at a location known to Genevieve.
    It contains the analytical work supporting this book. It is not necessary
    for the book to stand on its own, but it answers questions the book raises.
\end{enumerate}

\section*{What the Author Would Say}
\label{appendix:what-the-author-would-say}

``I don't know which possibility is true. I never did. I built this book to
work under all three because that's the honest position. Publish it if you
think it helps. The physics is coming regardless. The only question is whether
people are prepared.''
